\section{环把两种运算放进同一结构}
\label{sec:04-Discrete-Mathematics/08-Rings-and-Lattices/01}

要把用户 ID 摊到 6 个分片上，最省事的写法是乘一个常数再取模。取 \(h(k)=5k\bmod 6\)，把 \(0,1,2,3,4,5\) 依次送进去，得到 \(0,5,4,3,2,1\)——每个分片正好接到一个，顺序还被打乱了，正是想要的效果。把常数改成 2，同一行代码给出 \(0,2,4,0,2,4\)：三个分片承担了全部流量，另外三个一条记录也收不到。

差别不在 5 比 2 大。5 在模 6 意义下可以撤销，因为 \(5\times5=25\equiv1\)，乘 5 之后再乘一次 5 就回到原处；2 撤销不了。更麻烦的是 2 还带着另一条性质：

\[
2\times3\equiv0\pmod 6,
\]

两个非零的东西乘出了零。乘法把 6 个值压成 3 个，根源就在这里。

前面\hyperref[sec:04-Discrete-Mathematics/06-Groups-and-Symmetry/01]{``对称操作引出群''}给过一套处理运算的语言，但它在这里使不上劲：群一次只管一种运算。模 6 的加法确实是个规规矩矩的群，每个元素都有加法逆元，出问题的是乘法，而群的结构里根本没有第二种运算的位置，也就没有地方安放 \(a(b+c)=ab+ac\)——这条恰恰是把加法和乘法绑在一起的规律。环就是为了腾出这个位置。

\subsection{加法给足，乘法给得很少}

环 \(R\) 的加法那一半是照搬来的：封闭、结合、交换，有零元，每个元素有相反数，完整的交换群。乘法那一半要求低得多，只需封闭、结合，本书还默认存在单位元 1。两半之间靠分配律 \(a(b+c)=ab+ac\) 与 \((a+b)c=ac+bc\) 通信，除此之外乘法什么也不欠。

吝啬是有原因的。整数里除了 \(\pm1\)，没有任何元素有乘法逆元；方阵环 \(M_n(\R)\) 在 \(n>1\) 时连交换律都不成立，\(AB\) 和 \(BA\) 通常不是一回事。要是把``非零元可逆''或者``乘法交换''也写进定义，整数、多项式、矩阵这些最常打交道的对象会被一次性扫地出门。定义留得松，是为了让结论覆盖得广。

代价是每往上收紧一条要求，就换回一样能力，而下面这些例子分别停在不同的层上。

\begin{center}
\begin{tikzpicture}[x=1cm,y=1cm,font=\footnotesize,>=stealth]
  \node[font=\small,black!75] at (-0.2,4.55) {一种运算};
  \draw[black!55,fill=blue!8] (1.2,4.1) rectangle (3.1,4.9);
  \node at (2.15,4.5) {半群};
  \draw[black!55,fill=blue!12] (5.0,4.1) rectangle (6.9,4.9);
  \node at (5.95,4.5) {幺半群};
  \draw[black!55,fill=blue!18] (8.8,4.1) rectangle (10.7,4.9);
  \node at (9.75,4.5) {群};
  \draw[->] (3.2,4.5) -- (4.9,4.5);
  \node[align=center,black!70] at (4.05,5.25) {加单位元};
  \node[align=center,black!60] at (4.05,3.75) {有``不动''的\\ 那个元素};
  \draw[->] (7.0,4.5) -- (8.7,4.5);
  \node[align=center,black!70] at (7.85,5.25) {加逆元};
  \node[align=center,black!60] at (7.85,3.75) {\(ax=b\) 可解\\ 且解唯一};

  \node[font=\small,black!75] at (-0.2,1.35) {两种运算};
  \draw[black!55,fill=orange!8] (1.2,0.9) rectangle (2.9,1.7);
  \node at (2.05,1.3) {环};
  \draw[black!55,fill=orange!12] (4.3,0.9) rectangle (6.0,1.7);
  \node at (5.15,1.3) {交换环};
  \draw[black!55,fill=orange!18] (7.4,0.9) rectangle (9.1,1.7);
  \node at (8.25,1.3) {整环};
  \draw[black!55,fill=orange!26] (10.5,0.9) rectangle (12.2,1.7);
  \node at (11.35,1.3) {域};
  \draw[->] (3.0,1.3) -- (4.2,1.3);
  \node[align=center,black!70] at (3.6,2.05) {乘法交换};
  \draw[->] (6.1,1.3) -- (7.3,1.3);
  \node[align=center,black!70] at (6.7,2.05) {无零因子};
  \node[align=center,black!60] at (6.7,0.2) {两边同除\\ 变得安全};
  \draw[->] (9.2,1.3) -- (10.4,1.3);
  \node[align=center,black!70] at (9.8,2.05) {非零元可逆};
  \node[align=center,black!60] at (9.8,0.2) {线性方程组\\ 可以消元};
  \node[align=center,black!60,font=\scriptsize] at (2.05,0.1) {\(M_n(\R)\)};
  \node[align=center,black!60,font=\scriptsize] at (5.15,0.1) {\(\Z/6\Z\)};
  \node[align=center,black!60,font=\scriptsize] at (8.25,0.1) {\(\Z,\ F[x]\)};
  \node[align=center,black!60,font=\scriptsize] at (11.35,0.1) {\(\Q,\R,\ \Z/p\Z\)};
\end{tikzpicture}

\emph{同一件事做两遍：每加一条要求，可用的对象变少，可做的操作变多。上排是群那条线，只涉及一种运算；下排是环这条线。开头的分片函数活在 \(\Z/6\Z\) 里，它卡在``交换环''这一格，再往右一步就被 \(2\times3=0\) 挡住。}
\end{center}

图里最要紧的是右下两格。整数为什么可以放心地从 \(ab=ac\)、\(a\neq0\) 推出 \(b=c\)？不是因为整数特别，而是因为它落在``整环''那一格；\(\Z/6\Z\) 落在左边一格，这条推理在那里就不成立了。

\subsection{零因子让``两边同除''失效}

把模 6 的乘法整张摊开，问题一眼看得见。

\begin{center}
\begin{tikzpicture}[x=0.86cm,y=0.86cm,font=\small]
  \foreach \i in {0,...,5}{
    \foreach \j in {0,...,5}{
      \pgfmathtruncatemacro{\p}{mod(\i*\j,6)}
      \ifnum\p=0
        \ifnum\i>0
          \ifnum\j>0
            \fill[red!25] (\j-0.5,-\i-0.5) rectangle (\j+0.5,-\i+0.5);
          \fi
        \fi
      \fi
      \node at (\j,-\i) {\p};
    }
  }
  \foreach \k in {0,...,5}{
    \node[black!70] at (\k,0.95) {\k};
    \node[black!70] at (-0.95,-\k) {\k};
  }
  \draw[black!45] (-0.5,0.5) -- (5.5,0.5);
  \draw[black!45] (-0.5,-5.5) -- (5.5,-5.5);
  \draw[black!45] (-0.5,0.5) -- (-0.5,-5.5);
  \draw[black!45] (5.5,0.5) -- (5.5,-5.5);
  \foreach \k in {0,...,4}{
    \draw[black!25] (-0.5,-\k-0.5) -- (5.5,-\k-0.5);
    \draw[black!25] (\k+0.5,0.5) -- (\k+0.5,-5.5);
  }
  \draw[blue!70,thick] (-0.5,-4.5) rectangle (5.5,-5.5);
  \node[right,blue!70,font=\footnotesize] at (5.7,-5.0) {乘 5：整行是一个置换};
  \node[right,red!70,font=\footnotesize] at (5.7,-2.0) {乘 2：只落在 \(\{0,2,4\}\) 上};
\end{tikzpicture}

\emph{\(\Z/6\Z\) 的乘法表。红格是两个非零元素乘出零的位置，它们全部集中在 2、3、4 这几行——这三个数与 6 有公因子。第 5 行没有红格，且六个数各出现一次，所以乘 5 可以撤销，分片是均匀的。}
\end{center}

红格叫零因子：非零，但与某个非零元素相乘得零。有零因子在场，``两边同除''就不再可靠。从 \(ab=ac\) 只能得到 \(a(b-c)=0\)，接下来那一步``所以 \(b=c\)''其实调用了一条没写出来的假设——非零元素相乘不出零。方程 \(2x\equiv4\pmod6\) 有 \(x=2\) 和 \(x=5\) 两个解，而 \(5x\equiv4\pmod 6\) 只有一个。这两行代码看起来同构，行为不同。

没有零因子的交换环叫整环，每个非零元素都可逆的交换环叫域。\(\Z/n\Z\) 里的判据很干脆：\(a\) 可逆当且仅当 \(\gcd(a,n)=1\)，这一点\hyperref[sec:04-Discrete-Mathematics/04-Number-Theory/02]{``模运算与模逆''}用扩展 Euclid 算法给过构造。于是 \(n\) 取素数时所有非零元素都可逆，\(\Z/p\Z\) 是域；\(n\) 取合数时，它的每个真因子都是零因子。

\begin{proposition}
有限整环必定是域。
\end{proposition}

理由只有一行：固定非零元 \(a\)，映射 \(x\mapsto ax\) 在无零因子时是单射，有限集上的单射必是满射，于是 1 落在像里，那个原像就是 \(a^{-1}\)。有限性在这里不可省——\(\Z\) 是整环而不是域，正因为它无限。这条命题解释了为什么有限域的清单如此规整：有限的``可约分''直接等价于``可求逆''，中间没有第三种情况。

工程上撞见这件事的频率不低。散列与随机数生成里那些形如 \(x\mapsto (cx)\bmod 2^{64}\) 的混淆步骤，常数 \(c\) 一律取奇数，因为 \(\Z/2^{64}\Z\) 中的可逆元恰好是奇数；取了偶数，映射立刻塌缩，低位比特被成片抹平。哈希表容量选素数也是同一个考虑的另一面。

\subsection{理想是能安全``模掉''的那种子集}

模运算的做法可以推广：挑一批元素，宣布它们等于零，剩下的按同余类算。麻烦在于结果必须与代表元的选择无关。设想把子集 \(I\) 模掉，用 \(a\) 和 \(a+i\) 代表同一个类（\(i\in I\)），用 \(b\) 和 \(b+j\) 代表另一个，那么

\[
(a+i)(b+j)=ab+(aj+ib+ij),
\]

括号里的东西必须仍在 \(I\) 中，否则两次计算会落进不同的类，乘法在商上根本没定义。注意 \(aj\) 里的 \(a\) 是环里任意元素，所以要求不只是``\(I\) 对加法封闭''，而是 \(I\) 被整个环从两侧吸收：\(r\in R,\ a\in I\) 时 \(ra\) 与 \(ar\) 都留在 \(I\) 里。满足这条的子集叫理想。理想不是为了优美才这么定义的，它恰好是让商上的乘法有定义的最低要求。

整数模 \(n\) 就是 \(\Z/(n)\)，其中 \((n)\) 是 \(n\) 的全部倍数。同样的操作搬到多项式上更有意思：在 \(F[x]/(p(x))\) 里把 \(p(x)\) 当零，所有次数不低于 \(\deg p\) 的项按这条关系压下去。若 \(p\) 在 \(F\) 上不可约，商环是一个域；若 \(p=gh\) 可以分解，那么 \(g\) 与 \(h\) 在商里都非零而乘积为零——零因子又回来了，和 \(2\times3\equiv0\) 是同一个现象换了层外衣。

这一步是有限域的标准造法。AES 的字节运算发生在 \(\mathbb F_2[x]/(x^8+x^4+x^3+x+1)\) 里，Reed--Solomon 码把消息当作某个有限域上的多项式取值（见\hyperref[sec:04-Discrete-Mathematics/05-Coding-and-Polynomial-Methods/01]{``有限域与多项式插值''}）。这些系统之所以能用``解线性方程组''这类熟悉的手法处理字节，就是因为脚下那块地是域而不是随便一个环。

于是读代数式的方式可以定下来了：看到分式、消元、约分或者求逆，先问一句这个环给了你什么。分母可逆吗？消去律成立吗？乘法交换吗？同一串符号在 \(\Q\) 上、在 \(\Z/6\Z\) 上、在矩阵上会给出三种不同的结局，而分歧全部来自图里那几格。

环这条路线是往集合上叠运算，每叠一条要求，换一种可以放心做的操作。下一节走的是相反方向：不添运算，而是从已有的次序里读出运算来。\hyperref[sec:01-Mathematical-Language/02-Sets-Relations-Functions/03]{``偏序与可比较性''}里那个只会回答``谁在谁下面''的关系，配上``取上确界''和``取下确界''两个动作之后，也能变成一套可以演算的代数。
